% !PRINTED 44
\weberpage{44}
%<<CONT

\begin{equation}
\begin{aligned}
0&=B_0^{(0)}B_0^{(n)}-B_0^{(1)}B_1^{(n)}+B_0^{(2)}B_2^{(n)}-\cdots\pm B_0^{(n)}B_n^{(n)},\\
0&=-B_1^{(1)}B_1^{(n)}+B_1^{(2)}B_2^{(n)}-\cdots\pm B_1^{(n)}B_n^{(n)},\\
&\cdots\\
\pm 1&=\pm B_n^{(n)}B_n^{(n)},
\end{aligned}
\tag{6}
\end{equation}
und dies Formelsystem lässt sich in umgekehrter Reihenfolge mit Benutzung der Relation $B_\nu^{(n)}=B_{n-\nu}^{(n)}$ auch so darstellen:
\begin{equation}
\begin{aligned}
1&=B_0^{(n)}B_0^{(n)},\\
0&=B_1^{(n)}B_0^{(n)}-B_0^{(n-1)}B_1^{(n)},\\
0&=B_2^{(n)}B_0^{(n)}-B_1^{(n-1)}B_1^{(n)}+B_0^{(n-2)}B_2^{(n)},\\
&\cdots\\
0&=B_n^{(n)}B_0^{(n)}-B_{n-1}^{(n-1)}B_1^{(n)}+B_{n-2}^{(n-2)}B_2^{(n)}-\cdots\pm B_0^{(0)}B_n^{(n)},
\end{aligned}
\tag{7}
\end{equation}
oder auch in zusammenfassender Bezeichnung:
\begin{equation}
 \sum_{\nu=0}^{\mu}(-1)^\nu B_{\mu-\nu}^{(n-\nu)}B_\nu^{(n)}
 =\begin{cases}
 1,&\mu=0,\\
 0,&\mu=1,2,\ldots,n.
 \end{cases}
\tag{8}
\end{equation}


\begin{center}
{\large \S.\ 9.}

\vspace{0.5em}
{\large\gesperrt{Interpolation.}}
\end{center}

Wir wollen von den zuletzt gewonnenen Formeln eine Anwendung machen auf eine Aufgabe aus der Theorie der Interpolation.

\gesperrt{Es soll eine ganze rationale Function $n$ten Grades bestimmt werden, die für $n+1$ gegebene Werthe der Veränderlichen $x$ vorgeschriebene Werthe hat.}

Es handelt sich also um die Bestimmung der $n+1$ Coëfficienten $a_0,a_1,\ldots,a_n$ in
\begin{equation}
 f(x)=a_0x^n+a_1x^{n-1}+\cdots+a_{n-1}x+a_n
\tag{1}
\end{equation}
aus den $n+1$ Gleichungen
\begin{equation}
 f(\alpha_0)=A_0,\quad f(\alpha_1)=A_1,\quad \ldots,\quad f(\alpha_n)=A_n,
\tag{2}
\end{equation}
wenn $\alpha_0,\alpha_1,\ldots,\alpha_n$ die gegebenen Werthe von $x$, und $A_0,A_1,\ldots,A_n$ die entsprechenden Werthe von $f(x)$ sind.

Die Gleichungen (2) sind ein System von Gleichungen ersten Grades für die Unbekannten $a_0,a_1,\ldots,a_n$, die sich im Allgemeinen durch Determinanten auflösen lassen, wie wir im nächsten Abschnitt sehen werden. Wir kommen später auf diese Aufgabe zurück,
%CONT>>
